\abstract{
The lunar magnetic environment is characterized by the absence of a global dipole and the presence of localized crustal magnetic anomalies. As humanity prepares to return to the Moon through the Artemis program, understanding this highly heterogeneous magnetic environment is crucial for biological systems and astronaut health. In this study, we map the lunar crustal magnetic field at candidate Artemis, Apollo, and Chang'e landing sites, combining satellite magnetometer data with a systematic review of magnetobiology literature. We find that while all lunar landing sites present a hypomagnetic field (HMF) environment relative to Earth's geomagnetic field, significant heterogeneity exists. Some sites exhibit fields $< \SI{1}{\nano\tesla}$, whereas others present localized fields up to $\sim\SI{20}{\nano\tesla}$ at \SI{30}{\kilo\meter} altitude. We synthesize existing evidence on the effects of HMF on plant growth, microbiome stability, and biological development, and discuss the implications for future lunar surface operations. Our findings provide a framework for integrating magnetic field heterogeneity into site selection and experimental design for upcoming lunar biological investigations.
}

\keywords{Hypomagnetic field, Lunar magnetism, Artemis program, Magnetobiology, Space environments}
